% !TEX root = ../main.tex

\chapter{符号与标记}
\label{app:symbols}

为便于阅读，本文涉及的主要英文缩写、解剖标签、评价指标与数学符号汇总如下。
同一符号在不同章节中若具有局部定义，以首次出现处的上下文说明为准。

\section{英文缩写与术语}
\begin{longtable}{p{0.24\linewidth}p{0.68\linewidth}}
  \toprule
  符号或缩写 & 含义 \\
  \midrule
  ACL & 前交叉韧带（Anterior Cruciate Ligament） \\
  MLKI & 多发韧带膝关节损伤（Multi-Ligament Knee Injury） \\
  MRI & 磁共振成像（Magnetic Resonance Imaging） \\
  FSE & 快速自旋回波序列（Fast Spin Echo） \\
  DESS & 双回波稳态序列（Double Echo Steady-State） \\
  BML & 骨髓病变（Bone Marrow Lesion）；本文统计分析中主要指FSE高信号候选区域 \\
  OA & 骨关节炎（Osteoarthritis） \\
  PTOA & 创伤后骨关节炎（Post-traumatic Osteoarthritis） \\
  OAI-ZIB & Osteoarthritis Initiative-Zuse Institute Berlin 数据集 \\
  DICOM & 医学数字成像和通信格式（Digital Imaging and Communications in Medicine） \\
  NIfTI & 神经影像信息技术倡议格式（Neuroimaging Informatics Technology Initiative） \\
  UDA & 无监督域适应（Unsupervised Domain Adaptation） \\
  DANN & 域对抗神经网络（Domain Adversarial Neural Network） \\
  GRL & 梯度反转层（Gradient Reversal Layer） \\
  CycleGAN & 循环一致生成对抗网络（Cycle-Consistent Generative Adversarial Network） \\
  U-Net & 医学图像分割中常用的编码器-解码器卷积网络结构 \\
  ResNet34 & 34层残差网络骨干结构 \\
  BN & 批归一化（Batch Normalization） \\
  IN & 实例归一化（Instance Normalization） \\
  ASPP & 空洞空间金字塔池化（Atrous Spatial Pyramid Pooling） \\
  DIP & 深度图像先验（Deep Image Prior） \\
  SVR & 切片到体积重建（Slice-to-Volume Reconstruction） \\
  N4ITK & N4偏置场校正算法 \\
  EDT & 欧氏距离变换（Euclidean Distance Transform） \\
  kd-tree & 用于近邻检索的k维树数据结构 \\
  IQR & 四分位距（Interquartile Range） \\
  CI & 置信区间（Confidence Interval） \\
  \bottomrule
\end{longtable}

\section{组织标签与量化指标}
\begin{longtable}{p{0.24\linewidth}p{0.68\linewidth}}
  \toprule
  符号或缩写 & 含义 \\
  \midrule
  FB & 股骨（Femoral Bone） \\
  FC & 股骨软骨（Femoral Cartilage） \\
  TB & 胫骨（Tibial Bone） \\
  TC & 胫骨软骨（Tibial Cartilage） \\
  DSC & Dice相似系数（Dice Similarity Coefficient），用于衡量分割掩膜重叠程度 \\
  ASD & 平均表面距离（Average Surface Distance） \\
  RSD & 根均方表面距离（Root Mean Square Surface Distance） \\
  MSD & 最大表面距离（Maximum Surface Distance） \\
  $p$ & 统计学检验中的显著性概率 \\
  $U$ & Mann--Whitney U 检验统计量 \\
  $t$ & Welch 双样本 t 检验统计量 \\
  $\beta$ & 多变量回归模型中的回归系数 \\
  \bottomrule
\end{longtable}

\section{主要数学符号}
\begin{longtable}{p{0.24\linewidth}p{0.68\linewidth}}
  \toprule
  符号 & 含义 \\
  \midrule
  $I_{\mathrm{fixed}}$ & 配准中的固定图像 \\
  $I_{\mathrm{moving}}$ & 配准中的浮动图像 \\
  $T_{\mathrm{rigid}}$ & 三维刚性变换 \\
  $\mathbf{R}$ & 刚性变换中的旋转矩阵 \\
  $\mathbf{t}$ & 刚性变换中的平移向量 \\
  $\mathbf{x}$, $\mathbf{x}^{\prime}$ & 变换前后的三维空间坐标 \\
  $H(\cdot)$ & 图像信息熵 \\
  $NMI$ & 归一化互信息（Normalized Mutual Information） \\
  $L_{\mathrm{seg}}$ & 分割损失 \\
  $L_{\mathrm{adv}}$ & 域对抗损失 \\
  $\lambda$ & 梯度反转层或损失项的权重系数 \\
  $\rho_t$ & 第 $t$ 次迭代的动态损失比率 \\
  $\tau_{\mathrm{lower}}$, $\tau_{\mathrm{upper}}$ & 动态冻结策略中的下限与上限阈值 \\
  $\theta_G$ & 特征提取器或生成器的网络参数 \\
  $\theta_D$ & 域判别器的网络参数 \\
  $\epsilon$ & 防止分母为零的极小常数 \\
  $\mu$, $\sigma$ & Z-score归一化中的均值与标准差 \\
  $z$ & Z-score归一化后的体素强度 \\
  \bottomrule
\end{longtable}
